% GENERATED FILE. Edit canonical lessons, then run npm run generate:books.
% canonical-source-hash: fnv1a64:96cb88e1df1b3af2
% canonical-lessons: FR-C180-preparer, FR-C180-presenter, FR-C180-preter, FR-C180-proposer, FR-C180-ramasser

\chapter{To prepare, To introduce, To lend, To propose, To pick up}
\label{ch:fr-to-prepare-to-introduce-to-lend-to-propose-to-pick-up}
\hlchaptermodality{180}

\begin{chapteropening}
\textbf{By the end of this chapter:} \emph{I can say to prepare, to introduce, to lend, to propose and to pick up.}

Complete the last lesson of chapter 180: I can say to prepare, to introduce, to lend, to propose and to pick up.
\end{chapteropening}

\section[préparer]{préparer — to prepare}

\label{lesson:FR-C180-preparer}



\begin{quote}
Before the new one: say the French for to throw, then the French for to turn.
\end{quote}

\subsection*{You'll want to know: préparer}
\textbf{préparer} — \textquotedblleft{}to prepare\textquotedblright{}.

\begin{cousinweb}
\textbf{préparer} $\leftarrow$ Latin \emph{praeparare}, \textquotedblleft{}to make ready beforehand\textquotedblright{} — a learned borrowing. (EtymDB 2.1)
\end{cousinweb}

\begin{grammarlens}[title={What you've built}]
One more word for this chapter.
\end{grammarlens}

\subsection*{Your turn}
\begin{itemize}
\raggedright
  \item \emph{Say it:} \emph{préparer}
  \item \emph{Say it:} \emph{préparer}, once more
  \item \emph{Say it:} say \emph{préparer}
  \item \emph{Recall:} say the French for to throw, then the French for to turn, then say \emph{préparer} again
\end{itemize}

\begin{tcolorbox}[breakable,colback=teal!4,colframe=teal!35!black,title={Before you move on}]
What does préparer mean? (\textquotedblleft{}To prepare\textquotedblright{}.) Say it once more.
\end{tcolorbox}
\section[présenter]{présenter — to introduce}

\label{lesson:FR-C180-presenter}



\begin{quote}
Before the new one: say the French for to turn, then the French for to prepare.
\end{quote}

\subsection*{You'll want to know: présenter}
\textbf{présenter} — \textquotedblleft{}to introduce\textquotedblright{}.

\begin{cousinweb}
\textbf{présenter} $\leftarrow$ Latin \emph{praesentare}, \textquotedblleft{}to show, to place before\textquotedblright{} — a learned borrowing. (EtymDB 2.1)
\end{cousinweb}

\begin{grammarlens}[title={What you've built}]
One more word for this chapter.
\end{grammarlens}

\subsection*{Your turn}
\begin{itemize}
\raggedright
  \item \emph{Say it:} \emph{présenter}
  \item \emph{Say it:} \emph{présenter}, once more
  \item \emph{Say it:} say \emph{présenter}
  \item \emph{Recall:} say the French for to turn, then the French for to prepare, then say \emph{présenter} again
\end{itemize}

\begin{tcolorbox}[breakable,colback=teal!4,colframe=teal!35!black,title={Before you move on}]
What does présenter mean? (\textquotedblleft{}To introduce\textquotedblright{}.) Say it once more.
\end{tcolorbox}
\section[prêter]{prêter — to lend}

\label{lesson:FR-C180-preter}



\begin{quote}
Before the new one: say the French for to prepare, then the French for to introduce.
\end{quote}

\subsection*{You'll want to know: prêter}
\textbf{prêter} — \textquotedblleft{}to lend\textquotedblright{}.

\begin{cousinweb}
\textbf{prêter} $\leftarrow$ Middle French \emph{prester} $\leftarrow$ Latin \emph{praestare}, \textquotedblleft{}to provide, to vouch for\textquotedblright{} (Lewis \& Short) — inherited; the circumflex marks the lost \emph{s}.
\end{cousinweb}

\begin{grammarlens}[title={What you've built}]
One more word for this chapter.
\end{grammarlens}

\subsection*{Your turn}
\begin{itemize}
\raggedright
  \item \emph{Say it:} \emph{prêter}
  \item \emph{Say it:} \emph{prêter}, once more
  \item \emph{Say it:} say \emph{prêter}
  \item \emph{Recall:} say the French for to prepare, then the French for to introduce, then say \emph{prêter} again
\end{itemize}

\begin{tcolorbox}[breakable,colback=teal!4,colframe=teal!35!black,title={Before you move on}]
What does prêter mean? (\textquotedblleft{}To lend\textquotedblright{}.) Say it once more.
\end{tcolorbox}
\section[proposer]{proposer — to propose}

\label{lesson:FR-C180-proposer}



\begin{quote}
Before the new one: say the French for to introduce, then the French for to lend.
\end{quote}

\subsection*{You'll want to know: proposer}
\textbf{proposer} — \textquotedblleft{}to propose\textquotedblright{}.

\begin{cousinweb}
\textbf{proposer} $\leftarrow$ Old French \emph{proposer}, after Latin \emph{proponere}, \textquotedblleft{}to put forward\textquotedblright{} — a learned borrowing. (EtymDB 2.1)
\end{cousinweb}

\begin{grammarlens}[title={What you've built}]
One more word for this chapter.
\end{grammarlens}

\subsection*{Your turn}
\begin{itemize}
\raggedright
  \item \emph{Say it:} \emph{proposer}
  \item \emph{Say it:} \emph{proposer}, once more
  \item \emph{Say it:} say \emph{proposer}
  \item \emph{Recall:} say the French for to introduce, then the French for to lend, then say \emph{proposer} again
\end{itemize}

\begin{tcolorbox}[breakable,colback=teal!4,colframe=teal!35!black,title={Before you move on}]
What does proposer mean? (\textquotedblleft{}To propose\textquotedblright{}.) Say it once more.
\end{tcolorbox}
\section[ramasser]{ramasser — to pick up}

\label{lesson:FR-C180-ramasser}



\begin{quote}
Before the new one: say the French for to lend, then the French for to propose.
\end{quote}

\subsection*{You'll want to know: ramasser}
\textbf{ramasser} — \textquotedblleft{}to pick up\textquotedblright{}.

\begin{cousinweb}
\textbf{ramasser} — built from \emph{re-} and the verb for \textquotedblleft{}to heap up\textquotedblright{}, from the word for \textquotedblleft{}a mass\textquotedblright{}, $\leftarrow$ Latin \emph{massa}, \textquotedblleft{}a lump\textquotedblright{} (Lewis \& Short).
\end{cousinweb}

\begin{grammarlens}[title={What you've built}]
One more word for this chapter.
\end{grammarlens}

\subsection*{Your turn}
\begin{itemize}
\raggedright
  \item \emph{Say it:} \emph{ramasser}
  \item \emph{Say it:} \emph{ramasser}, once more
  \item \emph{Say it:} say \emph{ramasser}
  \item \emph{Recall:} say the French for to lend, then the French for to propose, then say \emph{ramasser} again
\end{itemize}

\begin{tcolorbox}[breakable,colback=teal!4,colframe=teal!35!black,title={Before you move on}]
What does ramasser mean? (\textquotedblleft{}To pick up\textquotedblright{}.) Say it once more.
\end{tcolorbox}
